\documentclass[10pt,a4paper]{amsart}
\usepackage{amsmath,amssymb,amsthm,graphicx}
\usepackage{verbatim}
\newenvironment{reference}{\comment}{\endcomment}
%\newenvironment{reference}{}{}
\newenvironment{slogan}{\comment}{\endcomment}
\newenvironment{history}{\comment}{\endcomment}

% For commutative diagrams we use Xy-pic
\usepackage[all]{xy}

% We use 2cell for 2-commutative diagrams.
\xyoption{2cell}
% Theorem environments.
%
\theoremstyle{plain}
\newtheorem{theorem}[subsection]{Theorem}
\newtheorem{proposition}[subsection]{Proposition}
\newtheorem{lemma}[subsection]{Lemma}

\theoremstyle{definition}
\newtheorem{definition}[subsection]{Definition}
\newtheorem{example}[subsection]{Example}
\newtheorem{exercise}[subsection]{Exercise}
\newtheorem{situation}[subsection]{Situation}

\theoremstyle{remark}
\newtheorem{remark}[subsection]{Remark}
\newtheorem{remarks}[subsection]{Remarks}

\numberwithin{equation}{subsection}

% Macros
%
\def\lim{\mathop{\mathrm{lim}}\nolimits}
\def\colim{\mathop{\mathrm{colim}}\nolimits}
\def\Spec{\mathop{\mathrm{Spec}}}
\def\Hom{\mathop{\mathrm{Hom}}\nolimits}
\def\Ext{\mathop{\mathrm{Ext}}\nolimits}
\def\SheafHom{\mathop{\mathcal{H}\!\mathit{om}}\nolimits}
\def\SheafExt{\mathop{\mathcal{E}\!\mathit{xt}}\nolimits}
\def\Sch{\mathit{Sch}}
\def\Mor{\mathop{\mathrm{Mor}}\nolimits}
\def\Ob{\mathop{\mathrm{Ob}}\nolimits}
\def\Sh{\mathop{\mathit{Sh}}\nolimits}
\def\NL{\mathop{N\!L}\nolimits}
\def\CH{\mathop{\mathrm{CH}}\nolimits}
\def\proetale{{pro\text{-}\acute{e}tale}}
\def\etale{{\acute{e}tale}}
\def\QCoh{\mathit{QCoh}}
\def\Ker{\mathop{\mathrm{Ker}}}
\def\Im{\mathop{\mathrm{Im}}}
\def\Coker{\mathop{\mathrm{Coker}}}
\def\Coim{\mathop{\mathrm{Coim}}}

% Boxtimes
%
\DeclareMathSymbol{\boxtimes}{\mathbin}{AMSa}{"02}

%
% Macros for moduli stacks/spaces
%
\def\QCohstack{\mathcal{QC}\!\mathit{oh}}
\def\Cohstack{\mathcal{C}\!\mathit{oh}}
\def\Spacesstack{\mathcal{S}\!\mathit{paces}}
\def\Quotfunctor{\mathrm{Quot}}
\def\Hilbfunctor{\mathrm{Hilb}}
\def\Curvesstack{\mathcal{C}\!\mathit{urves}}
\def\Polarizedstack{\mathcal{P}\!\mathit{olarized}}
\def\Complexesstack{\mathcal{C}\!\mathit{omplexes}}
% \Pic is the operator that assigns to X its picard group, usage \Pic(X)
% \Picardstack_{X/B} denotes the Picard stack of X over B
% \Picardfunctor_{X/B} denotes the Picard functor of X over B
\def\Pic{\mathop{\mathrm{Pic}}\nolimits}
\def\Picardstack{\mathcal{P}\!\mathit{ic}}
\def\Picardfunctor{\mathrm{Pic}}
\def\Deformationcategory{\mathcal{D}\!\mathit{ef}}
\usepackage{cleveref}
\crefname{originallemma}{original lemma}{original lemmas}
\Crefname{originallemma}{Original lemma}{Original lemmas}
\makeatletter
\newlabel{lemma-another-variant-local-criterion-flatness}{{00MO}{source}{}{sourcepointer-1}{}}
\newlabel{lemma-another-variant-local-criterion-flatness@cref}{{[originallemma][0][]00MO}{[1][0][]source}{}{}{}}
\makeatother
\crefname{originallemma}{original lemma}{original lemmas}
\Crefname{originallemma}{Original lemma}{Original lemmas}
\makeatletter
\newlabel{lemma-limit-module-essentially-finite-presentation}{{00QX}{source}{}{sourcepointer-2}{}}
\newlabel{lemma-limit-module-essentially-finite-presentation@cref}{{[originallemma][0][]00QX}{[1][0][]source}{}{}{}}
\makeatother
\crefname{originallemma}{original lemma}{original lemmas}
\Crefname{originallemma}{Original lemma}{Original lemmas}
\makeatletter
\newlabel{remark-Tor-ring-mod-ideal}{{00M6}{source}{}{sourcepointer-3}{}}
\newlabel{remark-Tor-ring-mod-ideal@cref}{{[originallemma][0][]00M6}{[1][0][]source}{}{}{}}
\makeatother
\title{Flatness of a colimit module is realised by one stage of the system}
\author{The Stacks Project Authors}
\date{stacks-project:algebra.tex:local-snapshot-20261012; GNU FDL 1.2 or later}
\begin{document}
\maketitle
\noindent\textit{Editorial note.} All delivered mathematical text is the author's own, copied from the pinned Stacks Project chapter file algebra.tex; no statement or proof step was rewritten, reordered or removed.\par\medskip
\noindent\textit{Editorial note.} Preamble provenance, as required for this source: the frozen snapshot stages the chapter's own preamble file beside it (preamble.tex, 217 lines, 8116 bytes, the same shared Stacks-wide preamble the published chow and etale-cohomology cards use); the first line of the pinned chapter file is \textbackslash{}input\{preamble\}, so no preamble text lives inside the chapter itself. The card takes that file's lines 8-18 (the verbatim package and the three \textbackslash{}comment-based environments reference/slogan/history at lines 9-12, section 13.2, together with the Xy-pic setup) and its lines 154 to the runtime end of the file (the native theorem declarations and mathematical macros). Not one character was changed, no package was added, no macro was defined, and the excerpt serves only as preamble, never as a span coordinate.\par\medskip
\noindent\textit{Editorial note.} The same-chapter pointers are located at their \textbackslash{}label lines inside the pinned chapter file; the cross-chapter pointers are located in the project's own tag list staged in the snapshot (sources/stacks-chow/tags.txt), at the verbatim lines carrying their canonical tags.\par\medskip
\noindent\textit{Editorial note.} The counted proof refers to three same-chapter statements by label (the setup lemma, a Tor remark and the local flatness criterion); they are recorded as navigation stubs pointed at their \textbackslash{}label lines, and the arguments consuming them stay inside the counted span.\par\medskip
\noindent\textit{Editorial note.} The counted spans contain no \textbackslash{}cite, so the card carries no bibliography.\par\medskip
\noindent\textit{Allowed standard prerequisites (editorial).}\begin{itemize}
\item The local criterion for flatness over a local ring. Invoked at the end of the counted proof (lines 32974-32976). Turns the vanishing of a Tor group over the residue field into flatness of the stage.
\item Commutation of tensor products with colimits. Used in the counted proof at lines 32959-32963. Moves the vanishing of the finitely many generators to a later stage of the system.
\end{itemize}\medskip
\ifdefined\hypertarget\hypertarget{sourcepointer-1}{}\fi
\noindent\textit{Source reference.} Original lemma 00MO: sources/stacks-algebra/pinned-original-source.tex, source lines 24138--24138. Navigation pointer for \textbackslash{}ref\{lemma-another-variant-local-criterion-flatness\} used inside the counted span (line 32977); the argument that consumes it is retained in the span.\par\medskip
\ifdefined\hypertarget\hypertarget{sourcepointer-2}{}\fi
\noindent\textit{Source reference.} Original lemma 00QX: sources/stacks-algebra/pinned-original-source.tex, source lines 32644--32644. Navigation pointer for \textbackslash{}ref\{lemma-limit-module-essentially-finite-presentation\} used inside the counted span (line 32942); the argument that consumes it is retained in the span.\par\medskip
\ifdefined\hypertarget\hypertarget{sourcepointer-3}{}\fi
\noindent\textit{Source reference.} Original lemma 00M6: sources/stacks-algebra/pinned-original-source.tex, source lines 18907--18907. Navigation pointer for \textbackslash{}ref\{remark-Tor-ring-mod-ideal\} used inside the counted span (line 32956); the argument that consumes it is retained in the span.\par\medskip
\begin{lemma}
\label{lemma-colimit-eventually-flat}
Let $R \to S$, $M$, $\Lambda$, $R_\lambda \to S_\lambda$, $M_\lambda$
be as in Lemma \ref{lemma-limit-module-essentially-finite-presentation}.
Assume that $M$ is flat over $R$.
Then for some $\lambda \in \Lambda$ the module
$M_\lambda$ is flat over $R_\lambda$.
\end{lemma}

\begin{proof}
Pick some $\lambda \in \Lambda$ and consider
$$
\text{Tor}_1^{R_\lambda}(M_\lambda, R_\lambda/\mathfrak m_\lambda)
=
\Ker(\mathfrak m_\lambda \otimes_{R_\lambda} M_\lambda
\to M_\lambda).
$$
See Remark \ref{remark-Tor-ring-mod-ideal}. The right hand side
shows that this is a finitely generated $S_\lambda$-module (because
$S_\lambda$ is Noetherian and the modules in question are finite).
Let $\xi_1, \ldots, \xi_n$ be generators.
Because $M$ is flat over $R$ we
have that $0 = \Ker(\mathfrak m_\lambda R \otimes_R M \to M)$.
Since $\otimes$ commutes with colimits we see there exists
a $\lambda' \geq \lambda$ such that each $\xi_i$ maps to
zero in
$\mathfrak m_{\lambda}R_{\lambda'} \otimes_{R_{\lambda'}} M_{\lambda'}$.
Hence we see that
$$
\text{Tor}_1^{R_\lambda}(M_\lambda, R_\lambda/\mathfrak m_\lambda)
\longrightarrow
\text{Tor}_1^{R_{\lambda'}}(M_{\lambda'},
R_{\lambda'}/\mathfrak m_{\lambda}R_{\lambda'})
$$
is zero. Note that
$M_\lambda \otimes_{R_\lambda} R_\lambda/\mathfrak m_\lambda$
is flat over $R_\lambda/\mathfrak m_\lambda$ because this last
ring is a field. Hence we may apply Lemma
\ref{lemma-another-variant-local-criterion-flatness}
to get that $M_{\lambda'}$ is flat over $R_{\lambda'}$.
\end{proof}

\end{document}
